\scanpage{17}
stable pour la stratification envisagée. De $=$, comme $VF=0$,
$V$ est nul sur $H'$, et \unsure{son} image étant contenue
dans $H^{10\,(p)}\inserted{=\operatorname{Ker}F}$\inserted{ (\unsure{donc $=$|étant $=$|équiv. à} $\operatorname{Ker}V=\operatorname{Im}F=H'$)}, \deleted{\unsure{$\operatorname{Ker}F(H^{(p)})$|car $FV=0$}},
et $=$ \unsure{égale|icelle} ; on trouve
que \uline{$H^{10\,(p)}$ est stable par \deleted{\unsure{connexion}}\inserted{la stratification}}. \deleted{donc}

\uline{Attention}, cela \uline{n'implique pas à priori que $H^{10}$ le
soit également}. Tout ce qu'on peut dire (valable
aussi pour $\Lambda$ une algèbre locale noeth. complète,
corps résiduel $k$ \ldots) c'est que si $S=\operatorname{Specf}(\Lambda)$,
et si on considère les deux images inverses
$E_1,E_2$ de $H$ sur \deleted{\unsure{$S$}} $T=\operatorname{Specf}(\Lambda\widehat\otimes\Lambda)=S\times S$,
\deleted{\unsure{alors}} et les deux images inverses $F_1,F_2$ de $H^{10}$,
et considérant l'isom de stratification
$u:E_1\xrightarrow{\sim}E_2$, d'où $u^{(p)}:E_1^{(p)}\mathbin{\unsure{-|\to}}E_2^{(p)}$,
alors $u^{(p)}(F_1^{(p)})=\deleted{\unsure{u^{(p)}}}\unsure{(|{}}F_2^{(p)}$, i.e.
$u(F_1)^{(p)}=F_2^{(p)}$. Si donc $u(F_1)=G$, $F_2=G'$, $E=E_2$
on aura $G^{(p)}=G'^{(p)}$ comme \deleted{\unsure{élémts}}\inserted{sections} de
de $\operatorname{Grass}(E)$ sur $T$, ou ce qui revient au
$=$, ces sections deviennent égales \deleted{par} \unsure{changmt|changt} de base
$\underset{\substack{\shortparallel\\T\ \text{puissance}\\\text{p.ième}}}{T'}\xrightarrow{\pi}T$. On ne peut en général en
conclure $G=G'$, car en général $\pi$ n'est pas un
épim. Mais si \inserted{$S$ donc} $T$ est formellt lisse sur $k$,
% archive sheet number: 16
